% ---- /md-to-pdf 共通プリアンブル（LuaLaTeX 用・全レイアウトに埋め込まれる）----
% pandoc は数式に amsmath の環境（aligned 等）を、チェックボックスに
% amssymb の \square を使う。原稿に数式が出た時点で落ちるので必ず読む。
\usepackage{amsmath,amssymb}
\usepackage{graphicx}
\usepackage{adjustbox}
\usepackage{xcolor}
\usepackage{enumitem}
\usepackage{fvextra}
% pandocの表出力（longtable）対応。calcは列幅計算の \real を供給する
\usepackage{longtable,booktabs,array}
\usepackage{calc}
\usepackage{newunicodechar}
%% 和文フォントに無い記号を、同義の和文記号へ寄せる
\newunicodechar{⚪}{○}
\newunicodechar{⚫}{●}
\usepackage[hidelinks]{hyperref}
% jlreqでは \@currentcounter が none のことがあり、hyperref+longtableが
% \refstepcounter{none} を呼んで「No counter 'none' defined」で落ちる。回避用に定義。
\makeatletter
\@ifundefined{c@none}{\newcounter{none}}{}
\makeatother
% 行末・行長調整で和欧混在時の改行位置と均等割りを改善
\usepackage{luatexja-adjust}
\ltjsetparameter{xkanjiskip={0.25\zw plus 0.15\zw minus 0.08\zw},
                 kanjiskip={0\zw plus 0.08\zw minus 0.08\zw}}

% 画像の相対パス（img/... 等）を入力mdのフォルダ基準で解決する
% （Obsidian vaultでは img/ がvaultルートにあることが多いため、親ディレクトリも探索する）
\graphicspath{__GRAPHICSPATH__}

% pandoc は \includegraphics に alt= キーを付けるので graphicx に無視させる
\makeatletter
\define@key{Gin}{alt}[]{}
\makeatother

% citeproc 引用マクロ（pandoc 3.x の \citeproc 形式に対応）
% \NewDocumentCommand はLaTeX 2020-10-01以降のカーネル機能。
% TeX Live 2019以前では未定義なので xparse を読み込んで補う。
\makeatletter
\@ifundefined{NewDocumentCommand}{\usepackage{xparse}}{}
\makeatother
\NewDocumentCommand\citeproctext{}{}
\NewDocumentCommand\citeproc{mm}{%
  \begingroup\def\citeproctext{#2}\cite{#1}\endgroup}
\makeatletter
 % 引用の行分割を許可し、\cite の括弧を抑止（citeproc書式に委ねる）
 \let\@cite@ofmt\@firstofone
 \def\@biblabel#1{}
 \def\@cite#1#2{{#1\if@tempswa , #2\fi}}
\makeatother
% citeproc の参考文献リスト（CSLReferences）用の定義（pandoc common.latex より）
\newlength{\cslhangindent}
\setlength{\cslhangindent}{1.5em}
\newlength{\csllabelwidth}
\setlength{\csllabelwidth}{3em}
\newenvironment{CSLReferences}[2] % #1 hanging-indent, #2 entry-spacing
 {\begin{list}{}{%
  \setlength{\itemindent}{0pt}
  \setlength{\leftmargin}{0pt}
  \setlength{\parsep}{0pt}
  \ifodd #1
   \setlength{\leftmargin}{\cslhangindent}
   \setlength{\itemindent}{-1\cslhangindent}
  \fi
  % 項目間は詰める（CSLの entry-spacing #2 は採用しない）
  \setlength{\itemsep}{0.2\baselineskip}}}
 {\end{list}}
\newcommand{\CSLBlock}[1]{\hfill\break\parbox[t]{\linewidth}{\strut\ignorespaces#1\strut}}
\newcommand{\CSLLeftMargin}[1]{\parbox[t]{\csllabelwidth}{\strut#1\strut}}
\newcommand{\CSLRightInline}[1]{\parbox[t]{\linewidth - \csllabelwidth}{\strut#1\strut}}
\newcommand{\CSLIndent}[1]{\hspace{\cslhangindent}#1}

%%% 学会クラス固有マクロのフォールバック %%%
% 投稿規定に沿って書いた原稿（\ecaption や \figref を含む）を、学会クラス以外の
% レイアウトでも下見できるようにする。学会クラス側で定義済みなら \providecommand は
% 何もしないので、ipsj などの本番ビルドには影響しない。
\providecommand{\ecaption}[1]{}          % 図表の英文見出しは落とす（和文見出しのみ残る）
\providecommand{\CaptionType}[1]{}
\providecommand{\figref}[1]{図~\ref{#1}}
\providecommand{\tabref}[1]{表~\ref{#1}}
\providecommand{\urlj}[1]{\texttt{#1}}
\providecommand{\refdatej}[1]{（参照 #1）}
\providecommand{\profile}[3]{\par\noindent\textbf{#2}\quad #3\par\smallskip}
\makeatletter
\@ifundefined{acknowledgment}{%
  \newenvironment{acknowledgment}{\par\medskip\noindent\textbf{謝辞}\par\nobreak}{\par\medskip}}{}
\@ifundefined{biography}{%
  \newenvironment{biography}{\par\medskip\noindent\textbf{著者紹介}\par\nobreak\small}{\par}}{}
\makeatother

% pandoc ヘルパ
\providecommand{\tightlist}{\setlength{\itemsep}{0pt}\setlength{\parskip}{0pt}}
\newcommand{\pandocbounded}[1]{\begin{center}%
  \adjustbox{max width=0.66\linewidth,max totalheight=0.5\textheight}{#1}\end{center}}

% ==ハイライト== → 折り返し可能な色付き太字（Obsidianの蛍光ペン相当）
\newcommand{\hl}[1]{\textbf{\textcolor[HTML]{B00000}{#1}}}

% Obsidianの箇条書きは6段以上ネストしうるので上限を拡張しつつ詰める
\setlistdepth{9}
\renewlist{itemize}{itemize}{9}
\renewlist{enumerate}{enumerate}{9}
\setlist{leftmargin=1.15em,itemsep=1pt,parsep=0pt,topsep=2pt,partopsep=0pt}
\setlist[itemize]{label=\textbullet}
\setlist[itemize,2]{label=\normalfont\bfseries--}
\setlist[itemize,3]{label=\textopenbullet}
\setlist[itemize,4]{label=\textperiodcentered}
\setlist[itemize,5]{label=\textasteriskcentered}
\setlist[itemize,6]{label=\textperiodcentered}
\setlist[itemize,7]{label=\textperiodcentered}
\setlist[itemize,8]{label=\textperiodcentered}
\setlist[itemize,9]{label=\textperiodcentered}

% コードブロック（XML等）：小フォントでトークン境界折り返し（語中で切らない）
\RecustomVerbatimEnvironment{verbatim}{Verbatim}%
  {fontsize=\scriptsize,breaklines=true,breakanywhere=false,breakautoindent=true,%
   breaksymbolleft=\raisebox{0.2ex}{\scriptsize\color{gray!70}$\hookrightarrow$},%
   breaksymbolindentleft=0.6em,xleftmargin=1.2em,xrightmargin=0.5em,%
   frame=leftline,framerule=0.4pt,rulecolor=\color{gray!50}}
